\chapter{包壳材料堆内行为模型补充关系式}
\label{app:cladding_supplementary_models}


% 为兼容原有子公式交叉引用：分组后保留旧标签，并统一指向所在关系组的公式编号。
\makeatletter
\providecommand{\appendixeqlabelalias}[1]{%
  \begingroup
  \edef\@currentlabel{\p@equation\theequation}%
  \label{#1}%
  \endgroup}
\makeatother

本附录汇总正文\ref{sec:cladding_incore_models}中未展开的材料成分、备选物性关系、经验参数和失效模型。不同关系式所依据的材料批次、冷加工量、热处理状态及辐照条件并不相同；模型使用时应优先保证材料状态和适用范围一致，而不应仅依据计算曲线的连续性进行外推。

\section{材料成分与状态}
\label{app:cladding_material_compositions}

\subsection{D9奥氏体不锈钢}

D9是在316系奥氏体不锈钢基础上发展形成的钛稳定化合金。钛与碳形成稳定碳化物，可改善高温组织稳定性及抗蠕变性能。D9-C1和高纯度D9-C1P的名义成分见表\ref{tab:appendix_d9_composition}\cite{hofmanMetallicFuelsHandbook2019}。表中数值均为质量百分数，Fe为余量。

\begin{table}[htbp]
	\centering
	\small
	\caption{D9-C1与D9-C1P的名义化学成分（wt.\%）}
	\label{tab:appendix_d9_composition}
	
	\begin{tabular*}{0.92\textwidth}{
			@{\extracolsep{\fill}}
			lcc lcc
			@{}
		}
		\toprule
		元素 & D9-C1 & D9-C1P & 元素 & D9-C1 & D9-C1P \\
		\midrule
		C  & 0.05 & 0.05 & Mn & 2.00 & 2.00 \\
		Si & 1.00 & 0.90 & P  & $\leq0.020$ & $\leq0.010$ \\
		S  & $\leq0.010$ & $\leq0.007$ & Cr & 13.5 & 13.5 \\
		Ni & 15.5 & 15.5 & Mo & 1.50 & 1.65 \\
		Cu & $\leq0.040$ & $\leq0.020$ & Nb & $\leq0.050$ & $\leq0.010$ \\
		Al & $\leq0.050$ & $\leq0.010$ & Ti & 0.25 & 0.25 \\
		As & $\leq0.030$ & $\leq0.030$ & Ta & $\leq0.020$ & $\leq0.010$ \\
		V  & $\leq0.020$ & $\leq0.020$ & Zr & -- & $\leq0.001$ \\
		Co & $\leq0.050$ & $\leq0.031$ & O  & -- & $\leq0.005$ \\
		B  & $\leq0.001$ & $\leq0.001$ & N  & $\leq0.010$ & $\leq0.005$ \\
		Fe & 余量 & 余量 & & & \\
		\bottomrule
	\end{tabular*}
\end{table}

\subsection{HT9铁素体/马氏体钢}

HT9为约12Cr--1Mo--W--V系回火马氏体钢，其UNS牌号为S42100。典型成分范围见表\ref{tab:appendix_ht9_composition}\cite{hofmanMetallicFuelsHandbook2019}。

\begin{table}[htbp]
	\centering
	\small
	\caption{HT9的典型化学成分范围（wt.\%）}
	\label{tab:appendix_ht9_composition}
	
	\begin{tabular*}{0.85\textwidth}{
			@{\extracolsep{\fill}}
			lc
			@{\hspace{1.2cm}}
			lc
			@{}
		}
		\toprule
		元素 & 含量 & 元素 & 含量 \\
		\midrule
		C  & 0.17--0.23 & Mn & 0.40--0.70 \\
		P  & $\leq0.015$ & S  & $\leq0.010$ \\
		Si & 0.20--0.30 & Ni & 0.30--0.80 \\
		Cr & 11.0--12.5 & Mo & 0.80--1.20 \\
		Nb & $\leq0.05$ & W  & 0.40--0.60 \\
		Al & $\leq0.050$ & V  & 0.25--0.35 \\
		Fe & 余量 & & \\
		\bottomrule
	\end{tabular*}
\end{table}

\section{热物性备选关系式}
\label{app:cladding_alternative_thermal_models}

\subsection{SS316热导率备选模型}

Matolich的NASA报告在$255.37$--$1255.37\ \mathrm{K}$范围内给出了分段热导率关系\cite{matolich1965thermal}：
\begin{equation}
	k_{316}^{\mathrm{NASA}}(T)
	=\begin{cases}
		3.155A_0(T/1.8)^{B_0},
		&255.37\leq T\leq699.816\ \mathrm{K}\\
		3.155\left(C_0+D_0T/1.8\right),
		&699.816<T\leq1255.37\ \mathrm{K}
	\end{cases}
	\label{eq:appendix_ss316_nasa_thermal_conductivity}
\end{equation}
其中，$A_0=0.52048$、$B_0=0.43367$、$C_0=4.553$、$D_0=5.507\times10^{-3}$，式\eqref{eq:appendix_ss316_nasa_thermal_conductivity}已换算为SI单位。

\subsection{独立密度拟合}

若未采用正文式\eqref{eq:clad_density_from_thermal_strain}的质量守恒更新，可使用下列线性密度拟合：
\begin{equation}
	\label{eq:appendix_cladding_density_fits}
	\begin{aligned}
		\rho_{\mathrm{D9}}(T_C)
		&=7.98-4.3\times10^{-4}T_C,
		&&300<T_C<800\ ^\circ\mathrm C\\
		\rho_{\mathrm{HT9}}(T_C)
		&=7.778-3.07\times10^{-4}T_C,
		&&0\leq T_C\leq800\ ^\circ\mathrm C
	\end{aligned}
\end{equation}
\appendixeqlabelalias{eq:appendix_d9_density_fit}
\appendixeqlabelalias{eq:appendix_ht9_density_fit}
式中，密度单位为$\mathrm{g\,cm^{-3}}$。HT9在高于约$800\ ^\circ\mathrm C$后接近相变区，不宜继续使用同一线性关系。

\subsection{SS316 Legacy IFR弹性关系}

为兼容早期快堆计算，部分程序仍保留下列SS316关系：
\begin{equation}
	\label{eq:appendix_ss316_legacy_elasticity}
	\begin{aligned}
		E_{316}^{\mathrm{legacy}}
		&=2.15946\times10^{11}-7.07727\times10^7T\\
		\nu_{316}^{\mathrm{legacy}}&=0.31
	\end{aligned}
\end{equation}
\appendixeqlabelalias{eq:appendix_ss316_legacy_youngs_modulus}
\appendixeqlabelalias{eq:appendix_ss316_legacy_poisson_ratio}
式中，$E$的单位为Pa，$T$的单位为K。由于该关系的原始来源已无法可靠追溯，正文不将其作为推荐模型。

\section{辐照肿胀参数}
\label{app:cladding_swelling_parameters}

\subsection{SS316}

正文式\eqref{eq:clad_void_swelling_inventory}中，SS316的温度参数为\cite{halesBISONTheoryManual2016,hofmanMetallicFuelsHandbook2019}
\begin{equation}
	\label{eq:appendix_ss316_swelling_parameters}
	\begin{aligned}
		\beta&=\frac{T-773}{100}\\
		R(T)&=\exp\!\left(
		0.497+0.795\beta-0.0948\beta^2
		+0.908\beta^3-1.49\beta^4
		\right)
		\\
		&\quad+\exp[-8(\beta-1.35)^2]\\
		\alpha&=0.75\\
		\tau(T)&=
		\begin{cases}
			6.58-0.566\beta,&T<848\ \mathrm{K}\\
			4.3105+2.46\beta,&T\geq848\ \mathrm{K}
		\end{cases}
	\end{aligned}
\end{equation}
\appendixeqlabelalias{eq:appendix_ss316_swelling_beta}
\appendixeqlabelalias{eq:appendix_ss316_swelling_R}
\appendixeqlabelalias{eq:appendix_ss316_swelling_alpha}
\appendixeqlabelalias{eq:appendix_ss316_swelling_tau}
采用增量更新时，空洞项的变化率为
\begin{equation}
	\dot S_0
	=\frac{0.01R(T)\dot\Phi}
	{1+\exp\{\alpha[\tau(T)-\Phi]\}},
	\qquad
	\dot\varepsilon_{v,316}^{\mathrm{sw}}
	=\frac{\dot S_0}{(1-S_0)^2}
	\label{eq:appendix_ss316_swelling_rate}
\end{equation}

对于特定20\%冷加工316钢数据，也可用幂律近似表示肿胀：
\begin{equation}
	\frac{\Delta V}{V_0}=A_{\mathrm{sw}}\Phi^n
	\label{eq:appendix_cw316_power_law_swelling}
\end{equation}
其中，$A_{\mathrm{sw}}$和$n$需根据材料炉批、温度和辐照装置确定。不同试验中$n$约为1.5--2.0，表明肿胀机制可能随气泡向空洞主导过程转变。

\subsection{D9}

D9的肿胀率、孕育期及固态密度修正为
\begin{equation}
	\label{eq:appendix_d9_swelling_parameters}
	\begin{aligned}
		R_{\mathrm{D9}}(T)
		&=2.76\exp[-1.4\times10^{-4}(T-773.15)^2]\\
		\alpha_{\mathrm{D9}}&=0.75,
		\qquad
		\tau_{\mathrm{D9}}=11.9\\
		D_{\mathrm{D9}}(T,\Phi)
		&=0.01(1-\exp[-30\Phi])
		(-1.7\times10^{-4}T+0.241)
	\end{aligned}
\end{equation}
\appendixeqlabelalias{eq:appendix_d9_swelling_R}
\appendixeqlabelalias{eq:appendix_d9_swelling_alpha_tau}
\appendixeqlabelalias{eq:appendix_d9_solid_state_density_change}
相应的增量形式为
\begin{equation}
	\dot\varepsilon_{v,\mathrm{D9}}^{\mathrm{sw}}
	=\frac{0.01R_{\mathrm{D9}}(T)\dot\Phi}
	{1+\exp[0.75(11.9-\Phi)]}
	-0.3\dot\Phi\exp(-30\Phi)
	(-1.7\times10^{-4}T+0.241)
	\label{eq:appendix_d9_swelling_rate}
\end{equation}

\subsection{HT9}

HT9的模型参数为
\begin{equation}
	\label{eq:appendix_ht9_swelling_parameters}
	\begin{aligned}
		R_{\mathrm{HT9}}(T)
		&=0.085\exp[-1.0\times10^{-4}(T-673)^2]\\
		\alpha_{\mathrm{HT9}}&=0.75,
		\qquad
		\tau_{\mathrm{HT9}}=14.2\\
		D_{\mathrm{HT9}}(\Phi)
		&=0.0015[1-\exp(-0.1\Phi)]
	\end{aligned}
\end{equation}
\appendixeqlabelalias{eq:appendix_ht9_swelling_R}
\appendixeqlabelalias{eq:appendix_ht9_swelling_alpha_tau}
\appendixeqlabelalias{eq:appendix_ht9_solid_state_density_change}
其增量形式为
\begin{equation}
	\dot\varepsilon_{v,\mathrm{HT9}}^{\mathrm{sw}}
	=\frac{0.01R_{\mathrm{HT9}}(T)\dot\Phi}
	{1+\exp[0.75(14.2-\Phi)]}
	+1.5\times10^{-4}\dot\Phi\exp(-0.1\Phi)
	\label{eq:appendix_ht9_swelling_rate}
\end{equation}

\section{屈服强度、极限强度与延性}
\label{app:cladding_strength_models}

\subsection{D9与HT9屈服强度的一般形式}

D9和HT9的$0.2\%$偏移屈服强度采用同一类热激活经验形式\cite{hofmanMetallicFuelsHandbook2019}：
\begin{equation}
	\ln\!\left(\frac{F\sigma^*}{\sigma_y}\right)
	=A_y\left[
	\frac{(\sigma^*/G)^M}
	{\dot\varepsilon\exp(Q/RT)}
	\right]^\lambda
	\label{eq:appendix_cladding_yield_strength_general}
\end{equation}
式中，$F$为辐照修正因子，$\sigma^*$为温度相关特征应力，$G$为剪切模量，$\dot\varepsilon$为应变率。未辐照条件下取$F=1$。基本拟合参数见表\ref{tab:appendix_cladding_yield_parameters}；辐照修正函数应按所采用的金属燃料手册或程序版本保持一致。

\begin{table}[htbp]
	\centering
	\caption{D9与HT9屈服强度模型的基本参数}
	\label{tab:appendix_cladding_yield_parameters}
	\begin{tabular}{
			@{\extracolsep{\fill}}
			lccccc
			@{}
		}
		\toprule
		材料 & $A_y$ & $M$ & $\lambda$ & $Q$（cal/mol） & 适用快中子注量 \\
		\midrule
		D9  & 30000 & 5.5 & 0.22  & 70000 & $\Phi\leq12$ \\
		HT9 & 20000 & 5.0 & 0.172 & 82000 & $\Phi\leq11$ \\
		\bottomrule
	\end{tabular}
\end{table}

相应的特征应力和剪切模量为
\begin{equation}
	\label{eq:appendix_cladding_characteristic_stress}
	\begin{aligned}
		\sigma^*_{\mathrm{D9}}
		&=450-240\tanh\!\left(\frac{T_C-760}{300}\right),
		&
		G_{\mathrm{D9}}&=78600-38.41T_C\\
		\sigma^*_{\mathrm{HT9}}
		&=430-190\tanh\!\left(\frac{T_C-640}{225}\right),
		&
		G_{\mathrm{HT9}}&=89640-53.78T_C
	\end{aligned}
\end{equation}
\appendixeqlabelalias{eq:appendix_d9_characteristic_stress}
\appendixeqlabelalias{eq:appendix_ht9_characteristic_stress}
式中，应力和模量单位均为MPa。

\subsection{极限拉伸强度}

D9和HT9的极限拉伸强度由屈服强度缩放得到：
\begin{equation}
	\label{eq:appendix_d9_ht9_ultimate_strength}
	\begin{aligned}
		\sigma_{u,\mathrm{D9}}
		&=\sigma_{y,\mathrm{D9}}
		\left[
		0.295\left(1-0.9\tanh\frac{\Delta T+400}{150}\right)
		+1.055
		\right]\\
		\sigma_{u,\mathrm{HT9}}
		&=\sigma_{y,\mathrm{HT9}}
		\left[
		1.1-0.1\tanh\!\left(\frac{\Delta T-200}{200}\right)
		\right]
	\end{aligned}
\end{equation}
\appendixeqlabelalias{eq:appendix_d9_ultimate_strength}
\appendixeqlabelalias{eq:appendix_ht9_ultimate_strength}
其中，$\Delta T=T_C-T_{i,C}$为测试温度与辐照温度之差。未辐照材料应采用相应手册规定的参考辐照温度处理，而不应任意令$\Delta T=0$。

\subsection{总伸长率}

D9和HT9的最小总伸长率分别为
\begin{equation}
	\label{eq:appendix_d9_ht9_total_elongation}
	\begin{aligned}
		\varepsilon_{t,\mathrm{D9}}(\%)
		&=-2.22+1.64\times10^{-2}T_i
		+2.17\times10^{-4}\exp\!\left(\frac{T_i}{100}\right)
		-6.3\times10^{-3}T_t\\
		\varepsilon_{t,\mathrm{HT9}}(\%)
		&=9.97+1.34\times10^{-3}
		\exp\!\left(\frac{T_t}{100}\right)
	\end{aligned}
\end{equation}
\appendixeqlabelalias{eq:appendix_d9_total_elongation}
\appendixeqlabelalias{eq:appendix_ht9_total_elongation}
D9关系主要适用于$640<T_i<1010\ \mathrm{K}$且$T_t=T_i$或$T_t=T_i+110\ \mathrm{K}$的试验路径；HT9关系的温度范围约为$273<T_t<973\ \mathrm{K}$。

\subsection{SS316硬化备选模型}

简单幂律各向同性硬化可写为
\begin{equation}
	\sigma_y(p)=\sigma_y^0+Kp^n
	\label{eq:appendix_ss316_power_hardening}
\end{equation}
其中，$\sigma_y^0$、$K$和$n$均应依据材料状态和温度标定。若需描述逐渐饱和的各向同性硬化，可采用Chaboche形式：
\begin{equation}
	\label{eq:appendix_ss316_chaboche_isotropic_hardening}
	\begin{aligned}
		\dot r&=b(Q-r)\dot p\\
		r(p)&=Q[1-\exp(-bp)]\\
		\sigma_y(p)&=\sigma_y^0+r(p)
	\end{aligned}
\end{equation}
\appendixeqlabelalias{eq:appendix_ss316_chaboche_rate}
\appendixeqlabelalias{eq:appendix_ss316_chaboche_integrated}
\appendixeqlabelalias{eq:appendix_ss316_chaboche_yield}
上述参数不可跨材料制备工艺直接移植；增材制造316L、退火态316和冷加工316的硬化响应可能存在显著差异。

\section{蠕变模型参数与补充关系}
\label{app:cladding_creep_parameters}

\subsection{D9参数}

正文式\eqref{eq:d9_creep_rate}所用参数为\cite{hofmanMetallicFuelsHandbook2019}
\begin{table}[htbp]
	\centering
	\caption{D9热蠕变与辐照蠕变模型参数}
	\label{tab:appendix_d9_creep_parameters}
	\begin{tabular}{
			@{\extracolsep{\fill}}
			lccccccc
			@{}
		}
		\toprule
		$A_1$ & $A_2$ & $A_3$ & $A_4$ & $A_5$ & $A_6$ & $A_7$ & $\Omega$ \\
		\midrule
		67 & 9461 & $1.0\times10^{-4}$ & 3955 & 27680 & $6.14\times10^{13}$ & 17111 & 9 \\
		\bottomrule
	\end{tabular}
\end{table}
温度修正函数为
\begin{equation}
	R_{\mathrm{cr,D9}}(T)
	=1.9\times10^{-4}
	\exp[-3.0\times10^{-5}(T-823)^2]
	\label{eq:appendix_d9_creep_temperature_factor}
\end{equation}
该模型主要用于$350$--$750\ ^\circ\mathrm C$、低应变率以及约$2\times10^{15}\ \mathrm{n\,cm^{-2}\,s^{-1}}$量级的快中子通量。

在$973\ \mathrm{K}$下，D9和20\%冷加工SS316的试验蠕变率可分别近似为\cite{lathaThermalCreepProperties2008}
\begin{equation}
	\label{eq:appendix_d9_ss316_973k_creep}
	\begin{aligned}
		\dot\varepsilon_{\mathrm{D9}}
		&=6\times10^{-30}\sigma^{9.2}\\
		\dot\varepsilon_{316}
		&=6\times10^{-28}\sigma^9
	\end{aligned}
\end{equation}
\appendixeqlabelalias{eq:appendix_d9_973k_creep}
\appendixeqlabelalias{eq:appendix_ss316_973k_creep}
式中，$\dot\varepsilon$的单位为$\mathrm{s^{-1}}$，$\sigma$的单位为MPa。两式仅适用于相应试验温度和材料状态，不应作为全温区模型。

\subsection{HT9一阶、二阶和三阶热蠕变}

HT9总热蠕变由初生、稳态和三级分量构成：
\begin{equation}
	\bar\varepsilon_{\mathrm T}
	=\bar\varepsilon_{\mathrm{TP}}
	+\bar\varepsilon_{\mathrm{TS}}
	+\bar\varepsilon_{\mathrm{TT}}
	\label{eq:appendix_ht9_total_thermal_creep}
\end{equation}
各分量为
\begin{equation}
	\label{eq:appendix_ht9_thermal_creep_components}
	\begin{aligned}
		\bar\varepsilon_{\mathrm{TP}}
		={}&\left[
		C_1e^{-Q_1/(RT)}\bar\sigma
		+C_2e^{-Q_2/(RT)}\bar\sigma^4
		+C_3e^{-Q_3/(RT)}\bar\sigma^{1/2}
		\right]
		[1-e^{-C_4t}]\\
		\bar\varepsilon_{\mathrm{TS}}
		={}&\left[
		C_5e^{-Q_4/(RT)}\bar\sigma^2
		+C_6e^{-Q_5/(RT)}\bar\sigma^5
		\right]t\\
		\bar\varepsilon_{\mathrm{TT}}
		={}&C_7e^{-Q_6/(RT)}\bar\sigma^{10}t^4
	\end{aligned}
\end{equation}
\appendixeqlabelalias{eq:appendix_ht9_primary_creep}
\appendixeqlabelalias{eq:appendix_ht9_secondary_creep}
\appendixeqlabelalias{eq:appendix_ht9_tertiary_creep}
参数见表\ref{tab:appendix_ht9_thermal_creep_parameters}。

\begin{table}[htbp]
	\centering
	\scriptsize
	\caption{HT9热蠕变模型参数}
	\label{tab:appendix_ht9_thermal_creep_parameters}
	\begin{tabular}{cccccccc}
		\toprule
		$C_1$ & $C_2$ & $C_3$ & $C_4$ & $C_5$ & $C_6$ & $C_7$ \\
		\midrule
		13.4 & $8.43\times10^{-3}$ & $4.08\times10^{18}$ & $1.6\times10^{-6}$ & $1.17\times10^9$ & $8.33\times10^9$ & $9.53\times10^{21}$ \\
		\midrule
		$Q_1$ & $Q_2$ & $Q_3$ & $Q_4$ & $Q_5$ & $Q_6$ & $R$ \\
		\midrule
		15027 & 26451 & 89167 & 83142 & 108276 & 282700 & 1.987 \\
		\bottomrule
	\end{tabular}
\end{table}
活化能与气体常数采用一致的cal/mol单位制；$t$以s计，$T$以K计，$\bar\sigma$以MPa计。模型主要适用于$350$--$750\ ^\circ\mathrm C$及$0$--$250\ \mathrm{MPa}$。

HT9辐照蠕变为
\begin{equation}
	\dot{\bar\varepsilon}_{\mathrm I,\mathrm{HT9}}
	=\left[
	1.83\times10^{-4}
	+2.59\times10^{14}
	\exp\!\left(-\frac{73000}{1.987T}\right)
	\right]
	\phi_{22}\bar\sigma^{1.3}
	\label{eq:appendix_ht9_irradiation_creep}
\end{equation}

\subsection{Monkman--Grant关系}

最小蠕变率与断裂时间可用Monkman--Grant及其修正形式关联。对于本文所列试验数据，有
\begin{equation}
	\label{eq:appendix_cladding_monkman_grant}
	\begin{aligned}
		\dot\varepsilon_{m,316}^{-0.66}t_r&=0.839,
		&
		\frac{\dot\varepsilon_{m,316}^{-0.97}t_r}{\varepsilon_t}&=0.306\\
		\dot\varepsilon_{m,\mathrm{D9}}^{-0.722}t_r&=0.276,
		&
		\frac{\dot\varepsilon_{m,\mathrm{D9}}^{-0.85}t_r}{\varepsilon_t}&=1.77
	\end{aligned}
\end{equation}
\appendixeqlabelalias{eq:appendix_ss316_monkman_grant}
\appendixeqlabelalias{eq:appendix_d9_monkman_grant}
上述经验关系继承相应试验的单位制，使用前必须与原始数据保持一致。

\section{断裂时间与瞬态失效模型}
\label{app:cladding_rupture_models}

\subsection{SS316蠕变断裂关系}

ORNL关系将断裂时间表示为\cite{brinkmanElevatedTemperatureMechanicalProperties2001}
\begin{equation}
	\log_{10}t_r
	=C_h-0.01312\sigma
	-2.552\log_{10}\sigma
	+\frac{20880}{T}
	\label{eq:appendix_ss316_ornl_rupture_time}
\end{equation}
其中，$t_r$以h计，$\sigma$以MPa计，$T$以K计；平均材料批次可取$C_h=-11.870$。对于高强度和低强度包络，可分别采用$-11.065$和$-12.674$。

316FR的区域化关系为\cite{OnizawaDevelopmentCreepProperty2019}
\begin{equation}
	(T_C+273.15)
	[\log_{10}t_r+C+ZS]
	=A_0+A_1\log_{10}\sigma
	+A_2(\log_{10}\sigma)^2
	\label{eq:appendix_316fr_rupture_time}
\end{equation}
其中，$Z=0$表示平均行为，$Z=1.64$和$2.33$可分别用于近似$5\%$和$1\%$失效概率。短期和长期参数应按原始文献划分，不可混用。

\subsection{D9稳态断裂时间}

D9-C1P采用
\begin{equation}
	\log_{10}t_r
	=A+\frac{B}{T}
	+\frac{C}{T}\log_{10}\sigma_\theta
	\label{eq:appendix_d9_c1p_rupture_time}
\end{equation}
其参数见表\ref{tab:appendix_d9_c1p_rupture_parameters}。

\begin{table}[htbp]
	\centering
	\caption{D9-C1P断裂时间模型参数}
	\label{tab:appendix_d9_c1p_rupture_parameters}
	\begin{tabular}{
			@{\extracolsep{\fill}}
			lrrrr
			@{}
		}
		\toprule
		应力区间 & $A$ & $B$ & $C$ & 标准估计误差 \\
		\midrule
		高应力 & $-16.638$ & 34054 & $-7113$ & 0.26 \\
		低应力 & $-12.637$ & 22219 & $-3591$ & 0.18 \\
		\bottomrule
	\end{tabular}
\end{table}
式中，$t_r$以h计，$T$以K计，$\sigma_\theta$以MPa计。建议范围为$770<T<1100\ \mathrm{K}$且$t_r\geq2\ \mathrm h$。若两个分区同时有效，取较小的断裂时间进行保守评估。

D9-C1采用正文式\eqref{eq:clad_generic_rupture_time}，参数见表\ref{tab:appendix_d9_c1_rupture_parameters}。

\begin{table}[htbp]
	\centering
	\caption{D9-C1断裂时间模型参数}
	\label{tab:appendix_d9_c1_rupture_parameters}
	\begin{tabular}{
			@{\extracolsep{\fill}}
			lrrrr
			@{}
		}
		\toprule
		应力区间 & $A$ & $B$ & $C$ & $D$ \\
		\midrule
		高应力 & $-10.727$ & 29732 & $-7414.8$ & 0 \\
		中应力 & $-53.034$ & 61542 & $-21246.8$ & 18.361 \\
		低应力 & $-17.141$ & 32605 & $-6803.5$ & 0 \\
		\bottomrule
	\end{tabular}
\end{table}
建议范围为$640<T<1100\ \mathrm{K}$且$t_r\geq1\ \mathrm h$。应力分区边界随温度变化，应依据金属燃料手册或程序实现逐区判断。

\subsection{D9瞬态失效关系}

D9短时瞬态模型综合考虑环向应力、温度、辐照温度、快中子注量和温升速率\cite{FanningSAS4ASASSYS1Safety2017}：
\begin{equation}
	\label{eq:appendix_d9_transient_rupture_time}
	\begin{aligned}
	\ln t_r
	={}&A+B\ln\!\left[
	\ln\!\left(\frac{\sigma^*}{\sigma_\theta}\right)
	\right]
	+\frac{Q}{T}
	+\mathrm{TRAMP}
	+C\tanh\Phi
	\\
	&+125\tanh^{10}\!\left(\frac{\sigma_\theta}{550}\right)
	\left[1-\tanh^{25}\!\left(\frac{T_i}{583}-F_0\right)\right]
	\tanh\!\left(\frac{\Phi}{2.5}\right)
	\end{aligned}
\end{equation}
其中
\begin{equation}
	\label{eq:appendix_d9_transient_auxiliary_functions}
	\begin{aligned}
		\mathrm{TRAMP}
		&=-0.28+1.18\tanh[-0.5\ln(\dot T)-1]\\
		\sigma^*
		={}&775-387.5
		\tanh\!\left(\frac{1000-T_i}{200}\right)
		\tanh\!\left(\frac{\Phi}{2}
		\right)
		\\
		&+125\tanh^{10}\!\left(\frac{\sigma_\theta}{550}\right)
		\left[1-\tanh^{25}\!\left(\frac{T_i}{583}-0.438\right)\right]
		\tanh\!\left(\frac{\Phi}{2.5}\right)
	\end{aligned}
\end{equation}
\appendixeqlabelalias{eq:appendix_d9_tramp}
\appendixeqlabelalias{eq:appendix_d9_transient_characteristic_stress}
模型参数为$A=-43.06$、$B=7.312$、$C=-1.73$和$Q=41339$。$t_r$以h计，$\sigma_\theta$以MPa计，$T$与$T_i$以K计，$\dot T$以$\mathrm{K\,s^{-1}}$计。

\subsection{HT9断裂时间}

金属燃料手册采用高、低应力两组参数：
\begin{equation}
	\log_{10}t_r
	=A+\frac{B}{T}
	+\frac{C}{T}\log_{10}\sigma_\theta
	\label{eq:appendix_ht9_mfh_rupture_time}
\end{equation}
高、低应力区的参数分别为
\begin{equation}
	\label{eq:appendix_ht9_mfh_rupture_parameters}
	\begin{aligned}
		(A,B,C)_{\mathrm H}
		&=(-32.490,\ 57781,\ -11800)\\
		(A,B,C)_{\mathrm L}
		&=(-35.173,\ 45858,\ -5563.1)
	\end{aligned}
\end{equation}
\appendixeqlabelalias{eq:appendix_ht9_mfh_high_stress_parameters}
\appendixeqlabelalias{eq:appendix_ht9_mfh_low_stress_parameters}
式中，下标$\mathrm H$和$\mathrm L$分别表示高应力区与低应力区；$t_r$以h计，$\sigma_\theta$以MPa计，适用温度约为$700$--$1100\ \mathrm{K}$。

WHC瞬态模型写为\cite{HuangTransientFailureBehavior1994}
\begin{equation}
	\label{eq:appendix_ht9_whc_failure_model}
	\begin{aligned}
		t_r&=\theta\exp\!\left(\frac{Q}{RT}\right)\\
		\ln\theta
		&=A_T+12.47\ln\!\left[
		\ln\!\left(\frac{730}{\sigma_\theta}\right)
		\right]\\
		A_T&=
		\begin{cases}
			24.942-0.153T_C
			+9.488\times10^{-5}T_C^2,
			&T<1144.15\ \mathrm{K}\\
			-36.1+1.5\tanh[(T-1144.15)/80],
			&1144.15\leq T<1393.15\ \mathrm{K}\\
			-36.1+1.5\tanh(179/80),
			&T\geq1393.15\ \mathrm{K}
		\end{cases}
	\end{aligned}
\end{equation}
\appendixeqlabelalias{eq:appendix_ht9_whc_rupture_time}
\appendixeqlabelalias{eq:appendix_ht9_whc_dorn_parameter}
\appendixeqlabelalias{eq:appendix_ht9_whc_temperature_factor}
其中，$T_C=T-273.15$，$Q=70.17\ \mathrm{kcal\,mol^{-1}}$，$t_r$以s计。该模型面向快速升温条件，不应与长期稳态关系无条件拼接。

\subsection{HT9受约束空洞长大模型}

受约束空洞长大模型将晶界空洞半径$a$的变化率表示为\cite{TVERGAARD1985447,KARAHAN2010283}
\begin{equation}
	\dot a
	=\frac{\dot V}{4\pi a^2h(\psi)},
	\qquad
	h(\psi)
	=\frac{1}{\sin\psi}
	\left[
	\frac{1}{1+\cos\psi}-\frac{\cos\psi}{2}
	\right]
	\label{eq:appendix_ht9_cavity_radius_rate}
\end{equation}
当空洞以$2b$为间距周期分布时，可将$a=b$作为局部贯通失效条件。空洞体积增长率由扩散和幂律蠕变贡献叠加：
\begin{equation}
	\dot V=\dot V_{\mathrm d}+\dot V_{\mathrm c}
	\label{eq:appendix_ht9_cavity_volume_rate}
\end{equation}
扩散贡献为
\begin{equation}
	\dot V_{\mathrm d}
	=4\pi D_{\mathrm{gb}}
	\frac{\sigma_n-(1-f)\sigma_s}
	{\ln(1/f)-(3-f)(1-f)/2}
	\label{eq:appendix_ht9_diffusional_cavity_growth}
\end{equation}
其中，$\sigma_s=2\gamma_s\sin\psi/a$，
$D_{\mathrm{gb}}=D_B\delta_B\Omega_a/(RT)$，且
\begin{equation}
	f=\max\!\left[
	\left(\frac{a}{b}\right)^2,
	\left(\frac{a}{a+1.5L}\right)^2
	\right],
	\qquad
	L=\left(\frac{D_{\mathrm{gb}}\sigma_e}
	{\dot\varepsilon_e^{\mathrm c}}\right)^{1/3}
	\label{eq:appendix_ht9_cavity_area_fraction}
\end{equation}
该模型对初始空洞半径和间距$b$高度敏感，必须由相应材料的显微组织或失效试验标定。


